The preceding chapter established that instability in RL poses a major barrier to reliably training LLMs for query augmentation tasks. To design experiments that address this issue, it is essential to understand how existing research has analyzed and mitigated instability in RL. This chapter reviews the body of work most relevant to that goal, focusing on three key perspectives. First, it examines studies that investigate how hyperparameter configurations, such as learning rate, clipping range, and regularization strength, affect the stability and convergence of policy-gradient methods, particularly PPO. Second, it explores how model architecture and scale influence the behavior of LLMs when used as policy makers within reinforcement learning frameworks. Finally, it reviews algorithmic approaches developed to improve stability, including TRPO, PPO and its recent variant, GRPO. Together, these strands of literature provide the theoretical and empirical foundations for the experimental design presented in the following chapter.

\subsection{Sensitivity to Hyperparameters}
To address the instability issues with PPO, several studies have analysed on the hyperparameters. \citeauthor{eimer2023hyperparameters}~\cite{eimer2023hyperparameters} conducted a comprehensive analysis of hyperparameter sensitivity across multiple reinforcement learning environments and algorithms, including PPO. Their study showed that seemingly minor changes in hyperparameter values can lead to large fluctuations in return and convergence behavior. In particular, they identified three hyperparameters as especially critical for PPO stability:
\begin{itemize}
    \item Learning rate, which directly affects the size of policy updates and can cause divergence when set too high,
    \item Clip range, which constrains the magnitude of policy updates and controls the trade-off between stability and learning progress, and
    \item Entropy coefficient, which influences exploration and prevents premature convergence but can destabilize learning when overemphasized.
\end{itemize}
Their results highlight that these parameters interact nonlinearly, making tuning difficult and environment dependent. Even within the same algorithm, the optimal configuration varies across tasks, suggesting that stability is not a fixed property of PPO but a function of its hyperparameter landscape.

Building on this line of work, \citeauthor{adkins2024method}~\cite{adkins2024method} introduced a general framework for measuring hyperparameter sensitivity in reinforcement learning. They proposed a systematic approach based on variance analysis to quantify how sensitive model performance is to small perturbations in each parameter. Their study emphasized not only the learning rate and clip range but also other influential parameters such as the discount factor, batch size, and number of epochs. By jointly analyzing these factors, they demonstrated that hyperparameter sensitivity often arises from complex interactions rather than single parameter effects. For instance, larger batch sizes can improve gradient stability but may require proportionally smaller learning rates or fewer epochs to maintain steady updates.

Together, these studies provide a foundation for understanding and addressing PPO instability. They suggest that investigating stability requires a systematic, controlled exploration of the hyperparameter space rather than isolated ablation studies. In this project, these findings guide the experimental design: focusing on the learning rate, clip range, KL or entropy coefficients, batch size, and number of epochs as the primary variables influencing training stability.

\subsection{Large Language Models}
To identify a suitable policy model for reinforcement learning in query augmentation, it is essential to review existing large language model families that balance accessibility, scale, and performance. The Qwen2 series~\cite{team2024qwen2}, developed by Alibaba Group, represents one of the strongest open-source options currently available. Released in mid-2024, Qwen2 spans a wide range of model sizes, from 0.5 billion to 72 billion parameters, which enables experiments to be conducted across multiple scales. Trained on a diverse multilingual corpus, Qwen2 models have demonstrated state-of-the-art performance among open-weight models on a variety of benchmarks, including reasoning, coding, and general knowledge. Their instruction-tuned variants also perform competitively on dialogue and retrieval tasks, suggesting strong generalization to text generation and query reformulation. All model checkpoints are publicly released on Hugging Face, making them practical for fine-tuning and reproducible experimentation. These qualities make Qwen2 a compelling foundation for studying reinforcement learning stability in large language models.

To complement Qwen2, the LLaMA 3.2 family developed by Meta provides another strong option for use as a policy model. The smallest text-only variants contain approximately one and three billion parameters, making them tractable for fine-tuning experiments~\cite{Llama3_2_3B}. The instruction-tuned version of LLaMA 3.2 is optimized for multilingual dialogue, information retrieval, and summarization tasks, demonstrating language modeling capabilities that are directly relevant to query augmentation. Because the models are officially released and hosted on Hugging Face under Meta’s community license, they are readily accessible for experimentation. Evaluation results indicate that LLaMA 3.2 performs well on benchmarks assessing reasoning, comprehension, and retrieval behaviors~\cite{MetaLlama32Blog}. These results suggest that LLaMA 3.2 has a solid baseline ability to generate semantically rich text, which is essential for effective query reformulation in retrieval settings.

Another promising candidate is the Mistral family, which consists of efficient and high performing open weight language models~\cite{jiang2023mistral7b}. Its publicly available variants include both dense and Mixture-of-Experts architectures within the 7B parameter range. Mistral models are designed to deliver strong performance while maintaining relatively low computational requirements. Benchmark results show that Mistral performs well on reasoning and general natural language processing tasks~\cite{jiang2023mistral7b}, placing it among the top models of comparable size in terms of general purpose understanding and generation. The models are available through major repositories such as Hugging Face, providing easy access for fine-tuning and controlled experimentation. Their efficiency allows greater flexibility for running multiple experiments or hyperparameter studies, making Mistral particularly suitable for analyzing the stability of PPO training. Its optimized design may also reveal different stability characteristics compared with Qwen2 or LLaMA due to differences in architecture and capacity.

Empirical evidence suggests that model size has a significant influence on both performance and stability during reinforcement learning fine-tuning. The study by ~\citeauthor{zhao2025echo}~\cite{zhao2025echo} shows that larger models tend to reinforce pretraining biases more strongly, while smaller models exhibit different amplification dynamics. This implies that model scale moderates how robust or volatile a policy becomes during reinforcement learning. Similarly, ~\citeauthor{huang2024n+}~\cite{huang2024n+} demonstrates that as model size increases, improvements in response quality also increase, yet the optimal hyperparameter configurations shift in complex ways. These findings indicate that stability and learning dynamics change nonlinearly with scale.

Together, these studies highlight the importance of including multiple model sizes in the experiments presented in this thesis. This approach allows not only the comparison of retrieval performance but also the examination of how training stability evolves with scale. The key question is whether larger models maintain smoother reward curves and more consistent optimization, or whether they experience greater instability under the same training conditions that work for smaller models. Understanding this relationship is crucial for identifying the most stable and reliable large language models for reinforcement learning in retrieval applications.

\subsection{RL Algorithms}
The stability of reinforcement learning is not determined solely by model architecture or hyperparameters but also by the choice of learning algorithm. Over the past decade, several policy gradient methods have been developed to improve stability, efficiency, and convergence reliability in high dimensional environments. Among these, algorithms such as TRPO and PPO have become foundational due to their principled approaches to constraining policy updates and preventing performance collapse. These algorithms address one of the central challenges in reinforcement learning, the tendency for large or poorly regularized parameter updates to destabilize training. This section reviews the key RL algorithms relevant to this study, focusing on their mathematical formulation, stability mechanisms, and suitability for fine-tuning large language models.

\subsubsection{Trust Region Policy Optimization}
Trust Region Policy Optimization, introduced by \citeauthor{schulman2015trust}~\cite{schulman2015trust}, is a foundational policy gradient algorithm designed to improve stability in RL by constraining how much the policy changes at each update. The key idea is to maximize a surrogate objective while enforcing a trust region constraint on the divergence between the new and old policies, thereby avoiding overly large, destabilizing updates.

Vanilla policy gradient methods can suffer from instability: a policy update that seems beneficial on sampled trajectories may lead to catastrophic performance degradation when applied broadly. TRPO addresses this by limiting each update to lie within a “safe” region, measured by Kullback–Leibler (KL) divergence, ensuring near-monotonic improvement under mild assumptions.

Let $\theta$ denote the policy parameterized by $\theta$. Define the advantage function $A^{\pi_{\text{old}}}(s, a)$ under the old policy. The TRPO update solves the following constrained optimization problem:
$$
\max_{\theta} L_{\pi_{\text{old}}}(\theta) = \mathbb{E}_{s \sim \rho_{\text{old}}, a \sim \pi_{\text{old}}} \left[ \frac{\pi_{\theta}(a \mid s)}{\pi_{\text{old}}(a \mid s)} A^{\pi_{\text{old}}}(s, a) \right]
$$
$$
\text{subject to} \quad \mathbb{E}_{s \sim \rho_{\text{old}}} \left[ \text{KL}(\pi_{\text{old}}(\cdot \mid s) \parallel \pi_{\theta}(\cdot \mid s)) \right] \leq \delta,
$$
where $\delta$ is a step size parameter controlling the maximum allowed change in policy per iteration. Here $\rho_{\text{old}}(s)$ is the discounted state visitation distribution under the old policy.

To make this tractable, TRPO linearizes the objective around $\theta_{\text{old}}$ and quadratically approximates the KL constraint, yielding a trust-region subproblem that can be solved via conjugate gradient methods. A backtracking line search is then used to ensure the updated policy satisfies the KL bound and improves the surrogate objective. In practice, the update takes the form:
$$\theta_{\text{new}} = \theta_{\text{old}} + \alpha \hat{x},$$
where $\hat{x}$ is the search direction computed by solving $H \hat{x} = g$ with $H$ being the Hessian (or Fisher information) matrix associated with the KL divergence and $g$ the policy gradient. The coefficient $\alpha$ is chosen via line search to satisfy the KL constraint.

\subsubsection{Proximal Policy Optimization}
PPO, introduced by \citeauthor{schulman2017proximal}~\cite{schulman2017proximal}, is one of the most widely adopted algorithms in reinforcement learning. It was designed as a simplified and more computationally efficient alternative to TRPO, while preserving its key principle: restricting policy updates to prevent large, destabilizing changes. PPO achieves this by introducing a clipped surrogate objective that penalizes excessive deviations between the new and old policies, thereby maintaining training stability without requiring second-order optimization or explicit trust-region constraints.

\begin{equation}
\begin{split}
    & \mathcal{L}^{PPO}(\theta) = \mathbb{E}[q \sim P(Q), \rho \sim \pi_{\theta_{old}}({O}|q)] \\ 
    & \frac{1}{|{o}|}\sum_{t=1}^{|{o}|} \min\left[\frac{\pi_{\theta}(o_t|q, {o}_{<t})}{\pi_{\theta_{old}}(o_t|q, {o}_{<t})}A_t, \text{clip}\left(\frac{\pi_{\theta}(o_t|q, {o}_{<t})}{\pi_{\theta_{old}}(o_t|q, {o}_{<t})}, 1-\epsilon, 1+\epsilon\right)A_t\right]
\end{split}
\end{equation} 

Here, ${\pi_{\theta}} \ $ and ${\pi_{\theta_{ref}}}$ denote the current and old policy models respectively, while $q$ and $o$ represent the input and generated output sampled from the dataset and the old policy. The hyperparameter $\epsilon$ specifies the clipping range that limits how far the probability ratio between the new and old policies can move in a single update. $ A_t $ denotes the advantage function at timestep $t$, computed from the observed reward and the estimated value function. The clipping mechanism ensures that the objective only increases when the policy improvement remains within a small, trusted region, effectively providing a soft constraint analogous to TRPO’s hard KL divergence limit.

Unlike TRPO, which relies on complex second-order updates and conjugate-gradient optimization, PPO uses straightforward stochastic gradient ascent, making it more scalable and easier to implement for large models such as LLMs. The algorithm alternates between collecting trajectories under the current policy and optimizing the clipped objective for several epochs on minibatches of this data. This iterative process improves the policy while maintaining stability, even in high dimensional action spaces.

PPO’s simplicity and robustness have made it the default choice for RL with LLMs. Its clipping objective acts as an implicit regularizer, preventing rapid divergence from the reference model and reducing the risk of reward over-optimization or policy collapse. In the context of this study, PPO serves as the baseline algorithm for analyzing training stability, with experiments designed to evaluate how hyperparameters and model size affect the evolution of KL divergence, entropy, and policy loss throughout the optimization process.

\subsubsection{Group Relative Policy Optimization}
GRPO is a variant of PPO introduced to improve learning stability by leveraging group-based feedback rather than single-sample rewards~\cite{shao2024deepseekmath}. While PPO evaluates each trajectory or output independently, GRPO generates a group of outputs for each state and uses the average behavior of that group to guide policy updates. This group-level averaging smooths out noise in the reward signal and reduces variance in the policy gradients, which can lead to more stable training.

The GRPO objective modifies the PPO clipped-loss form by sampling $G$ outputs $\left\{o_{i}\right\}_{i=1}^{G}$ from the old policy for each query, computing group-averaged advantages $\hat{A}_{i,t}$, and then optimizing:

\begin{equation}
\begin{split}
     & \mathcal{L}^{GRPO}(\theta) = \mathbb{E}[q \sim P(Q), \{o_i\}_{i=1}^G \sim \pi_{\theta_{old}}({O}|q)] \frac{1}{G}\sum_{i=1}^{G} \frac{1}{|{o}_i|}\sum_{t=1}^{|{o}_i|} \\
     & \left[\min\left[\frac{\pi_{\theta}(o_{i,t}|q, {o}_{i,<t})}{\pi_{\theta_{old}}(o_{i,t}|q, {o}_{i,<t})}\hat{A}_{i,t}, \text{clip}\left(\frac{\pi_{\theta}(o_{i,t}|q, {o}_{i,<t})}{\pi_{\theta_{old}}(o_{i,t}|q, {o}_{i,<t})}, 1-\epsilon, 1+\epsilon\right)\hat{A}_{i,t}\right] - \beta D_{KL}[\pi_{\theta}||\pi_{ref}]\right]
\end{split}     
\end{equation}

In this equation, $G$ is the group size, the number of outputs per query $\hat{A}_{i,t}$ is the advantage estimate for the $i$-th output at time $t$, $\beta$ is a coefficient scaling a KL penalty to prevent excessive divergence from a reference policy $\pi_{\text{ref}}$ and $\epsilon$ is the clipping threshold as in PPO.

The motivation behind GRPO is to trade off variance for bias in the policy gradient estimate. While group averaging introduces some bias, it significantly lowers the variance of updates, thus avoiding the erratic swings that plague standard PPO in high-dimensional or noisy settings. In practice, GRPO can lead to smoother reward curves, more stable KL divergence trajectories, and fewer failed policy updates.

Compared to PPO, GRPO needs to sample multiple outputs per query, which increases computational cost. However, this additional sampling is offset by more reliable gradient estimates and reduced need for extreme hyperparameter tuning to maintain stability. The addition of a KL penalty term further encourages cautious behavior and guards against over aggressive updates.

When applying RL to LLMs for tasks like query augmentation, reward signals are often noisy, delayed, and highly dependent on language model behavior. The group-based averaging strategy of GRPO is especially suitable in such contexts. It helps buffer the policy against outlier sequences or degenerate outputs. Since one of the research goals is to understand what algorithmic choices influence stability, evaluating GRPO alongside PPO offers a means to test whether group variance reduction leads to more consistent training, particularly across different model sizes and hyperparameter settings.